\section{General formalism}\label{sect_formalism}%
%

%
\subsection{Parton quasi-distributions}\label{subsect_definitions}%
%
Let us start with the following nucleon matrix element
\begin{align}
\label{def:ME}
 &\langle N(p)| \bar q (zv) \gamma_\alpha[zv,0] q(0) |N(p)\rangle  = 
\phantom{\frac{2v_\alpha}{v^2} (pv) \mathcal{I}^\parallel (v^2z^2\!,pvz,\mu)}
\notag\\
&\hspace*{2.0cm} = \frac{2v_\alpha}{v^2} (pv)\, \mathcal{I}^\parallel (v^2z^2\!,pvz,\mu)
\notag\\
&\hspace*{2.0cm}\quad+ 2\left(p_\alpha- \frac{v_\alpha}{v^2} (pv) \right) \mathcal{I}^\perp (v^2z^2\!,pvz,\mu)\,,
\end{align}
where
\begin{align}
\label{def:Wilson}
 [zv, 0] &= \text{Pexp}\Big[i g \int_0^z du\,  v^\mu A_\mu(uv)\Big]\,,
\end{align}
$p^\mu$ is the nucleon momentum, $p^2 = m^2$, $v^\mu$ is a given four-vector and $z$ a real number. All notations correspond to Minkowski space. In the following, we keep the normalization of the four-vector $v_\mu$ arbitrary, keeping in mind that $v^2 <0$ in the lattice calculation. We suppress flavor indices and tacitly assume considering the flavor-non-singlet combination of the matrix elements in what follows. We also neglect quark masses. 

The operator product in Eq.\,\eqref{def:ME} suffers from ultraviolet (UV) divergences and has to be renormalized. The argument $\mu$ of matrix elements $\mathcal{I}$ refers to the renormalization scale. The renormalization-scale dependence can be studied by going over to an effective theory~\cite{Craigie:1980qs,Dorn:1986dt,Ji:2017oey} such that the Wilson line is substituted by the propagator of an auxiliary field. For time-like separations the resulting theory is  the heavy-quark effective theory (HQET) and the renormalization factor of interest is the renormalization factor (squared) for the heavy-light quark current, see e.g.~\cite{Chetyrkin:2003vi}. We are not aware of a calculation for such a renormalization factor at space-like separations beyond one-loop order, however, to this accuracy there is no difference from the time-like case. 


The ``longitudinal'' and ``transverse'' invariant functions in Eq.~\eqref{def:ME} (with respect to $v_\mu$) correspond to particular projections of the matrix element that are employed in lattice calculations: 
\begin{align}\label{invfuncdef}
 \langle N(p)| \bar q (zv) \slashed{v} [zv,0]q(0) |N(p)\rangle  & = 2 (pv) \mathcal{I}^\parallel (v^2z^2,pvz,\mu)\,,
\notag\\
  \langle N(p)| \bar q (zv) \slashed{\epsilon}[zv,0] q(0) |N(p)\rangle &= 2 (p\epsilon) \mathcal{I}^\perp (v^2z^2,pvz,\mu)\,,
\end{align}
where $(\epsilon\cdot v) =0$. Eq.~(\ref{invfuncdef}) is the starting point for the construction of quasi-distributions.
 
The invariant functions  $\mathcal{I}^{\parallel}$, $\mathcal{I}^{\perp}$ \eqref{def:ME}  coincide at the tree level.  Assuming the power counting \begin{align}
\label{power-counting}
  z = \mathcal{O}(\eta)\,, \qquad p = \mathcal{O}(\eta^{-1})\,,\qquad \eta \to 0\,, 
\end{align} 
they can be written in terms of the position-space quark PDF~\cite{Ioffe:1969kf,Braun:1994jq,Hoyer:1996nr}:
\begin{align}
  \mathcal{I}^{\parallel(\perp)}(z^2v^2,pvz, \mu\sim 1/|vz|) &=  I(pvz, \mu_F \sim 1/|vz|)
%\int_{-1}^1 dx \, e^{ixz\,pv} q(x,\mu_F \sim 1/|z|)
\notag\\&\quad
 +\mathcal{O}(\alpha_s, \eta^2)\,, 
\end{align}
where  $\mu_F$ is the factorization scale, and 
\begin{align}
\label{ITD}
   I(pvz,\mu_F) &=  \int_{-1}^1 dx \, e^{ixz\,pv} q(x,\mu_F)\,.
\end{align} 
Here $q(x,\mu_F)$ for $x>0$ is the quark PDF, and for $x<0$ is the antiquark PDF. The position-space PDFs $I$ are known as the Ioffe-time distribution (ITD). To distinguish ITD $I$ from the functions $\mathcal{I}$, we refer to the functions $\mathcal{I}^{\parallel}$, $\mathcal{I}^{\perp}$  as the Ioffe-time \emph{quasi}-distributions (qITDs), following the terminology introduced in \cite{Ji:2013dva}.
 
The qPDFs $\mathcal{Q}^{\parallel},\mathcal{Q}^{\perp}$~\cite{Ji:2013dva} and the pPDF $\mathcal{P}$~\cite{Radyushkin:2017cyf}
are defined in terms of the qITDs by the Fourier transform
\begin{align}
\label{def:qPDF+pPDF}
\mathcal{Q}^{\parallel(\perp)}(x,p,\mu) &= |(pv)|\! \int\limits_{-\infty}^\infty \frac{d z }{2\pi} e^{-ix z (pv)} \, \mathcal{I}^{\parallel(\perp)}(z^2v^2, pvz,\mu )\,,
\notag\\
\mathcal{P}(x,z,\mu) &= |z| \int\limits_{-\infty}^\infty \frac{d (pv)}{2\pi} e^{-ix z (pv) } \,  \mathcal{I}^{\perp}(z^2v^2,pvz,\mu)\,.     
\end{align}
In what follows we will tacitly assume that $(pv)>0$ in the discussion of qPDFs and $z>0$ for pPDFs and drop the 
absolute value sign, for brevity.

In renormalization schemes with an explicit regularization scale, the Wilson line in Eq.\,\eqref{def:ME} suffers 
from an additional linear UV divergence that has to be removed. In dimensional regularization, this UV linear divergence reveals itself as a factorial growth of high orders of perturbative series~\cite{Beneke:1994sw}. The linear UV divergence can be removed by a mass renormalization associated with the Wilson line~\cite{Dotsenko:1979wb,Craigie:1980qs,Dorn:1986dt} or by a regularization-independent renormalization~\cite{Martinelli:1995vj,Stewart:2017tvs,Chen:2017mzz}. Given the multiplicative renormalizability of the quasi-PDF operator~\cite{Ji:2017oey,Ishikawa:2017faj,Green:2017xeu}, it can also be removed by considering a suitable ratio of matrix elements involving the same operator, e.g. by normalizing to the same matrix element at zero proton momentum~\cite{Orginos:2017kos}, 
\begin{align}
\label{n-qITD}
     \mathbf{I}(v^2z^2, pvz) =  \mathcal{I}(v^2z^2, pvz, \mu) / \mathcal{I}(v^2z^2, 0, \mu)\,, 
\end{align}
or, alternatively, to the vacuum expectation value
\begin{align}
\label{hat-qITD}
     \widehat{\mathbf{I}}(v^2z^2, pvz) =  \mathcal{I}(v^2z^2, pvz, \mu) / \mathcal{N}(v^2z^2,\mu)\,, 
\end{align}
where
\begin{align}
  \mathcal{N}(v^2z^2,\mu) =  \left(\frac{2 i N_c}{\pi^2 z^3 v^2}\right)^{-1} \langle 0| \bar q(zv) \slashed{v} [zv,0] q(0) |0\rangle\,.
\end{align}
Eqs.~(\ref{n-qITD}) and (\ref{hat-qITD}) will be our main focus in the present paper. By forming the ratio, the scale dependence cancels out (including the usual logarithmic renormalization) and one can define the scale-independent qPDF/pPDF
\begin{align}
\mathbf{Q}(x,p) &= |(pv)| \int_{-\infty}^\infty \frac{d z }{2\pi} e^{-ix z (pv)} \, \mathbf{I}(z, pv)\,,
\notag\\
\mathbf{P}(x,z) &= |z| \int_{-\infty}^\infty \frac{d (pv)}{2\pi} e^{-ix z (pv) } \,  \mathbf{I}(z,pv)\,,     
\end{align} 
and similarly for $\widehat{\mathbf{Q}}(x,p)$ and $\widehat{\mathbf{P}}(x,z)$.
The difference between these two options in the present context is that the normalization to the vacuum correlator 
does not affect the leading $\mathcal{O}(v^2z^2)$ power corrections that are subject of this work (since there is no gauge-invariant operator),
whereas the normalization to the value at zero momentum, as we will see, has a substantial effect. 
%Alternatively, the linear UV divergence can be tamed by  the nonperturbative subtractions, see e.g.~\cite{Martinelli:1995vj}.

%
\subsection{Light-ray OPE and target mass corrections}\label{subsect_OPE}%
%

The general approach to collinear factorization of QCD amplitudes in the position space is provided by the light-ray OPE~\cite{Anikin:1978tj,Anikin:1979kq,Balitsky:1987bk,Mueller:1998fv,Balitsky:1990ck,Geyer:1999uq}. For illustration consider the ``longitudinal'' projection. Specializing to the present case (forward matrix elements) we write
\begin{align}
\label{light-ray-OPE}
 \bar q(zv) \slashed{v} [zv,0] q(0) &= \int\limits_0^1\!d\alpha \, H^\parallel(z,\alpha,\mu_F)
\notag\\&\quad \times \Pi_{\rm l.t.}^{\mu_F}[\bar q (\alpha z v) \slashed{z} q(0)] + \ldots,
\end{align} 
where we include the renormalization factor in the coefficient function $H^\parallel = \delta(1-\alpha) + \mathcal{O}(\alpha_s)$ and set 
the renormalization scale $\mu$ to be equal to the factorization scale $\mu_F$. 
Finally,  $\Pi_{\rm l.t.}^{\mu_F}[\ldots]$ is the leading-twist projection operator defined below and
ellipses stand for the higher-twist contributions.

The leading-twist projection of the nonlocal quark-antiquark operator is 
defined as the generating function of \emph{renormalized} 
local leading-twist operators  (traceless and symmetrized over all indices)
\begin{align}
\Pi_{\rm l.t.}^{\mu_F}[\bar q (z v) \slashed{v} q(0)]
&= \sum_{n=1}^\infty \frac{z^{n-1}}{(n-1)!} v^{\mu_1}\ldots v^{\mu_n} O_{\mu_1\ldots\mu_n}^{n}(0)\,,
\label{eq:LTprojector}
\end{align}
where 
\begin{align}
  O_{\mu_1\ldots\mu_n}^{n}(z) &= 
\bar q(0)\gamma_{(\mu_1}\!\!\stackrel{\leftarrow}{D}_{\mu_2}\ldots \!\stackrel{\leftarrow}{D}_{\mu_{n})} q(0)\,.  
\end{align}
%
Here we indicate trace subtraction and symmetrization by enclosing the involved Lorentz indices in parentheses,
for example $O_{(\mu\nu)} = \frac12 (O_{\mu\nu}+ O_{\nu\mu}) - \frac14 g_{\mu\nu} O_{\lambda}^{~\lambda}$.

The light-ray OPE differs from the usual short-distance Wilson  expansion in local operators by imposing a different power counting. In the short-distance expansion one assumes that the distance between the quarks is small, $|z| \sim \eta \Lambda^{-1}_{\rm QCD}$ with $\eta\to 0$, and the operator matrix elements are of order unity in this limit, $\langle O_{\mu_1\ldots\mu_n}^{n} \rangle \sim \Lambda^{n}_{\rm QCD}$. In this case only a finite number of terms contribute to the r.h.s. of Eq.~\eqref{eq:LTprojector}, whereas the rest as well as the higher-twist operators must be added to 
higher orders of OPE, starting  from $\mathcal{O}(\eta^2)$. In other words, the relevant expansion parameter is the mass-dimension of an operator. 
The light-ray OPE assumes instead that the leading-twist operators scale as $\langle O_{\mu_1\ldots\mu_n}^{n} \rangle \sim \eta^{-n} \Lambda^{n}_{\rm QCD}$ so that $ z^n v^{\mu_1}\ldots v^{\mu_n} \langle O_{\mu_1\ldots\mu_n}^{n} \rangle  = \mathcal{O}(1)$. In this case, the series in~\eqref{eq:LTprojector} must be resummed to all orders, revealing its non-local "light-ray" structure. For a generic hadronic matrix element of leading twist operators one has
\begin{align}
 \langle p | O_{\mu_1\ldots\mu_n}^{n} |p\rangle \sim p_{(\mu_1}\ldots p_{\mu_n)} \langle\!\langle O^{n} \rangle\!\rangle\,, 
\end{align}      
where the reduced matrix element $\langle\!\langle O^{n} \rangle\!\rangle = \mathcal{O}(1)$. Therefore, the light-ray OPE is an adequate approximation if the hadron has large momentum, $|pv| = \mathcal{O}(\eta^{-1})$ and hence $ z\,pv =  \mathcal{O}(1)$. Higher-twist operators of the same dimension have smaller spin (by definition) and as a consequence, their matrix elements are power-suppressed \footnote{Strictly speaking, the expansion in powers of the large momentum corresponds to the classification in terms of the so-called collinear twist, see e.g.~\cite{Geyer:1999uq,Geyer:2000ig}.}.  

Note that the above power counting is applicable both in Minkowski and Euclidean space. In Minkowski space, one can 
go over to a different reference frame where all components of the momentum are of order $\Lambda_{\rm QCD}$ 
and simultaneously the separation between the quarks is almost light-like, $|v_\mu| = \mathcal{O}(1)$ but 
$v^2 = \mathcal{O}(\eta^2) \to 0$. 

On the calculation level, the light-ray OPE provides one with a convenient framework to operate with the leading-twist projected operators~\eqref{eq:LTprojector} avoiding the local expansion. Light-ray operators can be viewed as analytic operator functions of the separation 
between the quarks (all short-distance and light-cone singularities are subtracted). They satisfy the Laplace 
equation~\cite{Balitsky:1987bk}
\begin{align}
 \frac{\partial}{\partial v^\mu} \frac{\partial}{\partial v_\mu}  \Pi_{\rm l.t.}^{\mu_F}[\bar q (z v ) \slashed{v} q(0)] =0\,
\end{align}
with the boundary condition $\Pi_{\rm l.t.}^{\mu_F}\to 1 $ at $v^2 \to 0$. 
Explicit expressions for the projection operator $\Pi_{\rm l.t.}^{\mu_F}$ can be 
found in~\cite{Balitsky:1987bk,Geyer:1999uq,Geyer:2000ig,Braun:2011dg}. 
The light-ray OPE combined with the background field method is the standard technique, e.g., 
in light-cone sum rules~\cite{Balitsky:1989ry}, where it is used for the calculation of higher-twist contributions, 
and for the derivation of the evolution equations for off-forward parton distributions~\cite{Belitsky:1999hf,Belitsky:1998gc}. 
The renormalization group kernel for the evolution of flavor-singlet light-ray operators in the general off-forward kinematics 
is known to three-loop accuracy~\cite{Braun:2017cih}. 

The nucleon matrix element of the leading-twist projected operator~\eqref{eq:LTprojector} defines the leading-twist 
quark PDF,
\begin{eqnarray}
\lefteqn{ \langle N(p)| \Pi_{\rm l.t.}^{\mu_F}[\bar q (zv) \slashed{v} q(0)] |N(p)\rangle}
\notag\\ 
&=& 2 \int_{-1}^1\!dx\, \Pi_{\rm l.t.}[(p\cdot v) e^{ i x z (pv)}] q(x,\mu_F)\,, 
\end{eqnarray}
where~\cite{Balitsky:1990ck} 
\begin{align}
\label{Pi-exp}
\Pi_{\rm l.t.}[(pv) e^{ i x z(pv)}]
&= (pv)\Big[1 - \frac{1}{4} \frac{m^2 v^2}{(pv)^2} z\frac{d}{dz} \Big] e^{ ixz(pv)}
\nonumber\\&\quad
+ \mathcal{O}(m^4)\,. 
\end{align}
The second term in the square brackets is the leading nucleon mass correction,  which is the position-space counterpart of the Nachtmann target mass correction~\cite{Nachtmann:1973mr} in the deep-inelastic scattering (DIS). The all-order expression in powers of the nucleon mass for the leading-twist projected exponential function can be found in~\cite{Balitsky:1990ck}.

Taking the nucleon matrix element of the operator relation in Eq.~\eqref{light-ray-OPE} we obtain the 
factorization theorem for the ``longitudinal'' qITD in the form
\begin{align}
\mathcal{I}^\parallel   = &
 \int\limits_0^1 d\alpha \, H^\parallel(z,\alpha,\mu_F)
  \Big[1 - \frac{1}{4} \frac{m^2 v^2}{(pv)^2} z\frac{d}{dz} +\ldots \Big]\\ & \nonumber
\times I(\alpha z pv , \mu_F)
+\ldots\,,
\end{align}
where $I(\alpha z pv,\mu_F)$ is  ITD~\eqref{ITD} and ellipses stand for the higher-power target mass 
and ``genuine'' higher-twist corrections.

Making the Fourier transformation \eqref{def:qPDF+pPDF} one obtains, to the leading twist accuracy,
the factorization theorem for the ``longitudinal'' qPDF,
\begin{align}
\label{qPDF2}
 \mathcal{Q}^\parallel(x,p) &= \int_{-1}^1 \!\frac{dy}{|y|} C^\parallel_{\Qd}(\tfrac{x}{y},xp,\mu_F)
\Big[1 + \frac{1}{4} \frac{m^2 v^2}{(pv)^2} \frac{d}{dy} y +\ldots \Big] 
\notag\\&\quad 
\times  q(y,\mu_F)\,, 
\end{align}
where the coefficient function is given by 
\begin{align}
  C^\parallel_{\Qd}(\tfrac{x}{y},xp,\mu_F) &= (pv)|y|\!\int\limits_{-\infty}^{\infty}\!\frac{dz}{2\pi}\! \int\limits_0^1\!\! d\alpha\,
e^{i(pv)z(x-\alpha y)}
\notag\\&\quad \times H^\parallel(z,\alpha,\mu_F)\,. 
\end{align} 
In particular, at the leading order
\begin{align}
\label{target-parallel}
 \mathcal{Q}^\parallel(x,p) &= q(x) + \frac{1}{4} \frac{m^2 v^2}{(pv)^2} \big[x q'(x)  + q(x)\big]
+ \mathcal{O}\left({m^4}/{p^4}\right),
\end{align}
where $q'(x) = (d/dx) q(x)$. Note that the derivative applied to the quark PDF lowers the power $q(x) \stackrel{x \to 1}{\sim} (1-x)^p$ by one unit so that the mass correction is effectively enhanced by the factor $1/(1-x)$ as compared to the quark distribution itself at large Bjorken $x$. The similar enhancement of the target mass correction at $x\to 1$ is familiar from DIS~\cite{Nachtmann:1973mr}.

The target mass correction for the ``transverse'' qPDF can be calculated in a similar way,
starting from the nonlocal operator with an open Lorentz index \eqref{def:ME} and 
using the operator identity~\cite{Balitsky:1987bk}
\begin{align}
\Pi_{\rm l.t.}^{\mu_F}[\bar q (z v) \gamma_\alpha q(0)]
=  \int\limits_0^\infty \!dt\,\frac{\partial}{\partial v^\alpha} 
\Pi_{\rm l.t.}^{\mu_F}[\bar q ( t z v) \slashed{v} q(0)]\,,
\end{align}   
where the Wilson line between the quarks is implied. One obtains, at the tree level,
\begin{align}
\label{target-perp}
 \mathcal{Q}^\perp(x,p)  &=  q(x) + \frac{1}{4} \frac{ m^2 v^2 }{(pv)^2} \Big[ x q'(x) + 3 q(x) \Big]
\notag\\&\quad
-  \frac{1}{2} \frac{m^2 v^2}{(pv)^2} \theta(|x|<1)\!\! \int_{|x|}^1\! \frac{dy}{y} q(x/y) 
+ \mathcal{O}\left({m^4}/{p^4}\right),
\end{align}
and the target mass correction to the pPDF
\begin{align}
\label{target-pPDF}
 \mathcal{P}(x,z) &=  q(x) +  \frac14 z^2 v^2 m^2 x^2 \theta(|x|<1) \int_{|x|}^1 \frac{dy}{y} q(x/y)
\notag\\&\quad
+ \mathcal{O}\left({m^4}/{p^4}\right)\,.
\end{align}
The target mass corrections in Eqs.~(\ref{target-parallel}) and (\ref{target-perp}) agree with those derived in Ref.~\cite{Chen:2016utp} when expanded to $\mathcal O(m^2v^2/(pv)^2)$. Interestingly, the target mass correction to the pPDF is suppressed as $\mathcal{O}(1-x)$ at $x\to 1$ and not enhanced $\mathcal{O}(1/(1-x))$ 
in contrast to that to the qPDFs (\ref{target-parallel}) and the structure functions in DIS.

\subsection{Borel transform and renormalons}

The coefficient function $H^\parallel$ in Eq.~\eqref{light-ray-OPE} and the similar coefficient function 
$H^\perp$ in the $\MSbar$ scheme have the perturbative expansion 
\begin{align}
 H = \delta(1-\alpha) + \sum_{k=0}^\infty h_k a_s^{k+1}\,,\qquad a_s = \frac{\alpha_s(\mu)}{4\pi}\,, 
\end{align}
with factorially growing coefficients $h_k\sim k!$.

A convenient way to handle such a series is to consider the Borel transform
\begin{align}
  B[H] (w) = \sum_{k=0}^\infty \frac{h_k}{k!}\left(\frac{w}{\beta_0}\right)^k 
\end{align}
where powers of $\beta_0 = 11/3 N_C -2/3 n_f$ are inserted for the later convenience. The Borel image can be used as a generating function for the fixed-order coefficients
\begin{align}
 h_k = \beta_0^k \left(\frac{d}{dw}\right)^k  B[H] (w)\big|_{w=0}\,.
\end{align}
Moreover, the sum of the series can be obtained as the integral over positive values of the Borel parameter $w$
\begin{align}
\label{BorelIntegral}
H = \delta(1-\alpha)  + \frac{1}{\beta_0} \int_0^\infty \!dw\, e^{-w/(\beta_0 a_s)} B[H](w)\,. 
\end{align}
As it stands, the integral is not defined because the Borel transform generally has singularities on the integration path, known
as (infrared) renormalons. One can adopt a definition of the integral deforming the contour above or below the real
axis, or as the principle value. These definitions are arbitrary, and their difference, which is exponentially small in the 
coupling, must be viewed as an intrinsic uncertainty of perturbation theory that has to be removed by adding 
power-suppressed nonperturbative corrections. Another potential problem concerns the convergence of the Borel integral at 
$w\to\infty$. Since the quantity of interest depends on the single hard scale $1/|z v|$, dimension counting requires that the Borel
transform can be written as $(\mu^2 z^2 v^2 )^w$ times a function $F(w)$ of the Borel parameter and dimensionless kinematic variables.
Combining $(\mu^2 z^2 v^2)^w$ with $e^{-w/(\beta_0 a_s(\mu))}$ one sees that the (principal value) integral is convergent, provided 
the distance between the quarks $|zv|$ is sufficiently small compared to $1/\Lambda_{\rm QCD}$ and $F(w)$ does not increase exponentially
at $w\to\infty$, which is the case of all known examples. 

Naturally, a full all-order calculation cannot be performed. Instead, we employ the approximation~\cite{Beneke:1998ui,Beneke:2000kc}   
restricting ourselves to the perturbative series generated by the running-coupling effects in the one-loop diagrams, i.e. 
using QCD coupling at the scale of the gluon virtuality. 
Such contributions can be traced by computing  
the diagrams with the insertion of $k$ fermion loops in the one-loop diagram
and replacing  $- \frac23 n_f \mapsto   \beta_0 = \frac{11}{3} N_c - \frac23 n_f$, see Appendix A.
Another, equivalent technique \cite{Ball:1995ni} is based on the calculation of one-loop diagrams with an effective gluon mass.

%
\begin{figure*}[tbp]
\centering
 \includegraphics[width=0.95\textwidth]{BVZplots/bubblechain}
  %~\includegraphics[width=6.0cm]{BMJplots/fit-t4-qPDF}
\caption{Bubble-chain contribution to the coefficient function. The Wilson 
line factor is shown by the double dotted line.}
\label{fig:bubblechain}
\end{figure*}
% 

The singularity structure of the Borel transform can be extracted separately without explicit evaluation of the bubble-chain. It can be done by replacing the running coupling constant in the loop diagrams by its effective form,
\begin{align}
  \beta_0 a_s(-k^2) =  \int_0^\infty \!dw\, e^{\frac53 w} \left(\frac{\Lambda^2}{-k^2}\right)^w ,
\end{align}
where $\Lambda \equiv \Lambda_{\rm QCD}^{\MSbar}$ and the factor $e^{\frac53 w}$ represents the $\MSbar$ scheme. Such replacement leads to the modified form of the gluon propagator 
\begin{align}\label{modgluonprop}
  \frac{1}{-k^2-i\epsilon} ~\mapsto~ \frac{(\Lambda^2)^w}{(-k^2-i\epsilon)^{1+w}}\,.
\end{align}
At one-loop level, the pole structure in $w$ reproduces the pole structure for the Borel plane. In our work, we have used both methods of calculation for the cross-check of the result.


%
%%%%%%%%%%%%%%%%%%%%%%%%%%%%%%%%%%%%%%%%%%%%%%%%%%%%%
